\chapter{PENUTUP}
\label{chap:penutup}

% Ubah bagian-bagian berikut dengan isi dari penutup

\section{Kesimpulan}
\label{sec:kesimpulan}

Setelah dilakukan proses desain, simulasi, implementasi, pengujian, hingga
analisis dan pembahasan hasil, didapatkan kesimpulan sebagai berikut.

\begin{enumerate}[nolistsep]

  \item Sistem pengendalian tekanan pompa air berhasil dirancang dan
        direalisasikan menggunakan rangkaian daya PWM AC \emph{chopper}
        topologi diode \emph{full-bridge} dan dua IGBT dengan \emph{gate driver}
        HCPL-3120. Penentuan \emph{dead time} 1{,}5~$\mu$s terverifikasi
        pada pengukuran (1{,}48--1{,}49~$\mu$s) sehingga tidak terjadi
        \emph{shoot-through}. Rangkaian mampu mengatur tegangan RMS keluaran
        secara linear terhadap \emph{duty cycle} dengan selisih di bawah 1\%
        terhadap nilai teoritis, dan bentuk gelombang tegangan maupun arus
        hasil pengukuran sesuai dengan hasil simulasi PSIM.

  \item \emph{Neural network} berhasil diterapkan sebagai pengendali tekanan
        menggunakan arsitektur \emph{Multi-Layer Perceptron} 3-4-1 dengan
        masukan galat tekanan, perubahan galat, dan \emph{duty cycle}, serta
        keluaran berupa perubahan \emph{duty cycle} ($\Delta duty$). Model
        dilatih dari data operasi \emph{closed-loop} (\emph{behavior cloning})
        dan mencapai kinerja $R^2 = 0{,}915$ (RMSE $0{,}463$; MAE $0{,}329$)
        pada data uji, menunjukkan akurasi prediksi aksi kendali yang baik.

  \item Sistem kendali mampu menjaga tekanan air pada \emph{setpoint} dengan
        \emph{error} tunak $\leq$0{,}009~bar untuk variasi 1--4 keran terbuka.
        Pada pengujian perubahan \emph{setpoint} (0{,}25--0{,}35~bar),
        \emph{settling time} berkisar 4--7~detik. Pada pengujian rejeksi
        gangguan perubahan beban keran, kontroler memulihkan tekanan dalam
        9--15~detik. Keterbatasan terjadi pada kondisi seluruh keran tertutup
        dan pada gangguan pembukaan 4~keran serentak, ketika tekanan minimum
        pompa melampaui \emph{setpoint} sehingga aktuator mengalami saturasi
        (batas fisik pompa), bukan kegagalan kontroler.

\end{enumerate}

\section{Saran}
\label{chap:saran}

Untuk pengembangan penelitian lebih lanjut, disarankan hal-hal berikut.

\begin{enumerate}[nolistsep]

  \item Memperbaiki desain rangkaian \emph{snubber} agar mampu meredam
        lonjakan tegangan (\emph{spike}) dan osilasi (\emph{ringing}) secara
        lebih baik, serta menambahkan sistem pendingin (\emph{cooling}) pada
        \emph{snubber} dan komponen daya untuk menjaga kestabilan termal pada
        operasi kontinu.

  \item Mengintegrasikan sistem secara langsung dengan inverter (sumber DC/PLTS)
        sehingga menjadi satu kesatuan \emph{solar water pump inverter}, tidak
        hanya berperan sebagai pengatur tegangan AC.

  \item Memperbaiki \emph{settling time} kontroler NN agar respons sistem lebih
        cepat, baik pada pelacakan perubahan \emph{setpoint} (saat ini
        4--7~detik) maupun pada pemulihan gangguan perubahan beban keran (saat
        ini 9--15~detik).

  \item Melakukan analisa efisiensi daya dan kerugian daya pada rangkaian daya, serta menambahkan pengukuran arus beban untuk mengetahui
        karakteristik beban pompa air dan memprediksi konsumsi daya.

\end{enumerate}

